\documentclass{article}
\usepackage[T1]{fontenc}
\usepackage{lmodern}
\usepackage{graphicx}
\usepackage{amsmath,amssymb}
\usepackage[a4paper,margin=2cm]{geometry}
\usepackage[absolute,overlay]{textpos}
\setlength{\TPHorizModule}{1mm}
\setlength{\TPVertModule}{1mm}
\usepackage[indonesian]{babel}
\usepackage{hyperref}
\setlength{\emergencystretch}{2em}
% unit: Halaman 80

\begin{document}

% === Halaman 80 (llm) ===

Algoritma dekripsi kebalikan dari algoritma enkripsi. Langkah-langkahnya adalah sebagai berikut:

\begin{enumerate}
  \item Jika dua huruf terdapat pada baris bujursangkar yang sama maka tiap huruf diganti dengan huruf di kirinya.
  \item Jika dua huruf terdapat pada kolom bujursangkar yang sama maka tiap huruf diganti dengan huruf di atasnya.
  \item Jika dua huruf tidak pada baris yang sama atau kolom yang sama, maka huruf pertama diganti dengan huruf pada perpotongan baris huruf pertama dengan kolom huruf kedua. Huruf kedua diganti dengan huruf pada titik sudut keempat dari persegi panjang yang dibentuk dari tiga huruf yang digunakan sampai sejauh ini.
  \item Buanglah huruf X yang tidak mengandung makna.
\end{enumerate}

\end{document}
